\section{Multi-Agent Query Routing and Intent Classification}
\label{sec:agents}
\begin{figure}[!t]
\centering
\begin{tikzpicture}[
  font=\scriptsize, node distance=3.4mm,
  b/.style={draw, rounded corners=2pt, align=center, text width=5.7cm, minimum height=5.5mm, inner sep=2.5pt, fill=black!3},
  d/.style={draw, diamond, aspect=2.6, align=center, inner sep=0pt, text width=3.0cm, fill=black!5},
  a/.style={draw, rounded corners=2pt, align=center, text width=1.7cm, minimum height=5mm, inner sep=2pt, fill=black!8},
  pf/.style={-{Latex[length=1.7mm]}}
]
\node[b] (q) {query $q$ (already reformulated by context resolver)};
\node[d, below=of q] (multi) {conjunction split $\to$ $\ge 2$ different agents?};
\node[d, below=6mm of multi] (nav) {contains one of 23 navigation phrases?};
\node[d, below=6mm of nav] (room) {bare room-number regex matches?};
\node[d, below=6mm of room] (kw) {any domain keyword set non-empty?};
\node[d, below=6mm of kw] (rnn) {LSTM confidence $\ge 0.6$?};
\node[a, right=5mm of nav] (na1) {navigation\_agent};
\node[a, right=5mm of room] (na2) {navigation\_agent};
\node[a, right=5mm of kw] (best) {argmax domain agent};
\node[a, right=5mm of rnn] (map) {\code{INTENT\_TO\_AGENT}};
\node[a, below=6mm of rnn] (gen) {general\_agent};
\node[b, right=5mm of multi, text width=1.7cm] (mrun) {run each clause; concat};

\draw[pf] (q)--(multi);
\draw[pf] (multi)-- node[left]{no} (nav);
\draw[pf] (multi)-- node[above]{yes} (mrun);
\draw[pf] (nav)-- node[left]{no} (room);
\draw[pf] (nav)-- node[above]{yes} (na1);
\draw[pf] (room)-- node[left]{no} (kw);
\draw[pf] (room)-- node[above]{yes} (na2);
\draw[pf] (kw)-- node[left]{no} (rnn);
\draw[pf] (kw)-- node[above]{yes} (best);
\draw[pf] (rnn)-- node[above]{yes} (map);
\draw[pf] (rnn)-- node[left]{no} (gen);
\end{tikzpicture}
\caption{Supervisor routing cascade. The first four tests are deterministic
substring/regex matches; the learned LSTM classifier is reached only when
none of them fire. Every decision is logged with its trigger.}
\label{fig:routing}
\end{figure}

\begin{table}[!t]
\centering
\caption{Query-Routing Approaches and the Implemented Hybrid.}
\label{tab:routing}
\footnotesize
\renewcommand{\arraystretch}{1.2}
\begin{tabular}{@{}p{1.55cm}p{1.55cm}p{1.55cm}p{2.2cm}@{}}
\toprule
\textbf{Property} & \textbf{Rule / keyword} & \textbf{Learned classifier} & \textbf{LLM router} \\
\midrule
Latency & negligible & ms (LSTM) & seconds \\
Cost & none & training only & per-query inference \\
Auditability & full (literal match) & partial (weights) & low \\
Unseen phrasing & misses & generalizes (if trained well) & robust \\
Cold-start data need & none & labeled corpus & none \\
Determinism & total & total at inference & low \\
\midrule
\multicolumn{4}{@{}p{7.6cm}@{}}{\textbf{Implemented:} deterministic phrase/keyword cascade
\emph{first}; trained LSTM ($12$ classes) consulted only on no-match, accepted
only if $P(\text{intent})\ge 0.6$; \code{general\_agent} otherwise. Combines
rule auditability with a learned safety net for uncovered phrasings.} \\
\bottomrule
\end{tabular}
\end{table}


The assistant is organized as a supervisor and five specialist agents ---
\code{admission\_agent}, \code{academic\_agent}, \code{facilities\_agent},
\code{navigation\_agent}, \code{general\_agent}. Here ``multi-agent''
denotes an intent-driven routing architecture~\cite{wooldridge2009mas}, not
concurrency and not an LLM-driven agent loop~\cite{wu2023autogen}: the
supervisor makes no factual claims, calls no LLM, and merely classifies and
dispatches.

\subsection{Specialist Agents}
All five agents are thin wrappers over one shared function,
\code{run\_specialist}, which executes the identical pipeline of
Section~\ref{sec:rag} (hybrid retrieval $\to$ rerank $\to$ context
selection $\to$ confidence $\to$ gated generation $\to$ grounding check
$\to$ curated fallback). What differs between agents is \emph{which queries
the supervisor sends them}, plus two agent-specific behaviors:
\code{navigation\_agent} first attempts two deterministic lookup paths
(a panorama-signage spatial knowledge store built from real main-building
room numbers, and a PostgreSQL entity resolver) and only falls through to
RAG on a miss, forcing confidence to $1.0$/HIGH on an exact structured
hit; and \code{general\_agent} is a pure pass-through appropriate for a
no-clear-domain fallback. No agent has a private retriever or model client;
none answers from parametric memory. This shared-pipeline discipline is an
explicit architectural invariant.

\subsection{Routing Algorithm}
\code{supervisor.classify} (Algorithm~\ref{alg:routing},
Fig.~\ref{fig:routing}) is a cascade evaluated in a fixed order:

\begin{enumerate}
\item \textbf{Navigation phrases.} A $23$-entry list of high-specificity
substrings (``where is'', ``how do i get to'', ``which floor'', ``take me
to'', ``panorama'', \dots) is checked first, because these are unambiguous
location signals even when the object contains another domain's keyword
(``Where is the CSE \emph{department}?'' is a location question).
\item \textbf{Bare room number.} A regular expression matching an optional
``room''/``room no.''/``\#'' prefix followed by a $2$--$4$ digit number and
an optional letter routes ``room 202'' to navigation regardless of
surrounding verbs.
\item \textbf{Domain keyword scoring.} Each of four domain keyword sets is
scored by the count of its members occurring in $q$; the highest-scoring
non-empty set wins.
\item \textbf{LSTM fallback.} \emph{Only} if steps 1--3 all fail, the
trained LSTM classifier is consulted; if $P(\text{intent}) \ge 0.6$, its
predicted class is mapped to an agent via a $12\!\to\!5$ dictionary,
otherwise \code{general\_agent} is chosen.
\end{enumerate}
Every decision returns a human-readable reason, and the supervisor logs
\code{"Routing decision: query=\%r -> agent=\%s (\%s)"} at \code{INFO}
level --- an explicit functional requirement for demonstrability and
auditability. This is a deliberate design choice against LLM-router
approaches (Table~\ref{tab:routing}): it costs no routing latency, needs no
per-query model call, and any decision can be explained from a single log
line.

\begin{algorithm}[!t]
\caption{Multi-Agent Query Routing}
\label{alg:routing}
\begin{algorithmic}[1]
\Require query $q$; navigation phrases $\mathcal{N}$; keyword sets $\{\mathcal{K}_a\}$;
threshold $\tau = 0.6$
\Ensure agent name(s) and reason
\State \textbf{// multi-domain pre-check}
\State $\Gamma \gets \textsc{SplitOnConjunction}(q)$ \Comment{on `` and '' / `` ; ''}
\If{$|\Gamma| \ge 2$ \textbf{and} $|\{\textsc{Classify}(\gamma) : \gamma \in \Gamma\}| \ge 2$}
  \State \textbf{return} $\{(\gamma, \textsc{Classify}(\gamma)) : \gamma \in \Gamma\}$ \Comment{run each, concatenate}
\EndIf
\State $\hat q \gets \text{lowercase}(q)$
\ForAll{$\pi \in \mathcal{N}$}
  \If{$\pi \subseteq \hat q$}\ \textbf{return} (\code{navigation\_agent}, ``phrase '$\pi$' '') \EndIf
\EndFor
\If{\textsc{RoomNumberRegex}$(q)$}\ \textbf{return} (\code{navigation\_agent}, ``bare room number'') \EndIf
\ForAll{agent $a$}\ $\sigma_a \gets |\{w \in \mathcal{K}_a : w \subseteq \hat q\}|$ \EndFor
\State $a^\star \gets \arg\max_a \sigma_a$
\If{$\sigma_{a^\star} > 0$}\ \textbf{return} $(a^\star,\ \text{``keywords ''} \dots)$ \EndIf
\State $(\iota, p) \gets \textsc{ClassifyIntentLSTM}(q)$ \Comment{Alg.~\ref{alg:lstm}}
\If{$\iota \ne \varnothing$ \textbf{and} $p \ge \tau$}
  \State \textbf{return} $(\textsc{IntentToAgent}[\iota],\ \text{``LSTM } \iota\ (p)\text{''})$
\EndIf
\State \textbf{return} (\code{general\_agent}, ``no rule matched; LSTM not confident'')
\end{algorithmic}
\end{algorithm}

\subsection{Multi-Domain Queries}
A thin opt-in layer splits a query on an explicit conjunction (`` and '',
`` ; '') and classifies each clause with the same \code{classify}. Only if
two or more clauses land on genuinely \emph{different} agents is the
multi-domain path taken: each clause is answered independently by its agent
and the answers are concatenated with no second LLM merge call. A query
like ``What is the admission process and what is the fee structure?''
splits into two clauses that both classify to \code{admission\_agent} and
therefore falls straight through to the ordinary single-agent path.

\subsection{LSTM Intent Classifier}
The classifier (Algorithm~\ref{alg:lstm}) is a compact recurrent
network~\cite{hochreiter1997lstm}:
\begin{align}
\mathbf{x}_t &= \mathbf{E}\,\text{onehot}(w_t), \quad \mathbf{E}\in\mathbb{R}^{|\mathcal{V}|\times 32},\\
(\mathbf{h}_t,\mathbf{c}_t) &= \mathrm{LSTM}(\mathbf{x}_t, \mathbf{h}_{t-1}, \mathbf{c}_{t-1}), \quad \mathbf{h}_t\in\mathbb{R}^{32},\\
\mathbf{z} &= \mathbf{W}\,\mathrm{dropout}_{0.2}(\mathbf{h}_{T}) + \mathbf{b}, \quad \mathbf{W}\in\mathbb{R}^{12\times 32},\\
P(\text{class} = j \mid q) &= \mathrm{softmax}(\mathbf{z})_j,
\end{align}
where $w_1{:}w_T$ are the query's lowercase alphanumeric tokens,
$\mathbf{h}_T$ is the final hidden state after a packed sequence pass, and
the classifier head produces logits over $12$ classes. Vocabulary
$\mathcal{V}$ has $283$ entries (plus \code{<pad>}/\code{<unk>}); the same
minimal tokenizer as BM25 is used so that room numbers and acronyms are not
fragmented. Inference applies \code{softmax} and returns the arg-max class
with its probability; if artifacts are missing the classifier returns
$(\varnothing, 0)$ and the supervisor treats that as ``no opinion.''

\begin{algorithm}[!t]
\caption{LSTM Intent Classification Pipeline}
\label{alg:lstm}
\begin{algorithmic}[1]
\Require query $q$; vocab $\mathcal{V}$; trained params $(\mathbf{E},\text{LSTM},\mathbf{W},\mathbf{b})$; classes $\mathcal{I}$ ($|\mathcal{I}|{=}12$)
\Ensure $(\text{intent},\ \text{probability})$
\If{artifacts absent} \textbf{return} $(\varnothing,\ 0)$ \EndIf
\State $[w_1,\dots,w_T] \gets \textsc{Tokenize}(q)$ \Comment{\code{[a-z0-9]+}}
\If{$T = 0$} \textbf{return} $(\varnothing,\ 0)$ \EndIf
\State $\text{ids} \gets [\,\mathcal{V}.\text{get}(w_t, \langle\text{unk}\rangle)\,]_{t=1}^{T}$
\State $\mathbf{z} \gets \textsc{IntentLSTM}(\text{ids},\ \text{lengths}{=}T)$ \Comment{embed $\to$ pack $\to$ LSTM $\to$ dropout $\to$ linear}
\State $\mathbf{p} \gets \mathrm{softmax}(\mathbf{z})$;\ $j^\star \gets \arg\max_j p_j$
\State \textbf{return} $(\mathcal{I}[j^\star],\ p_{j^\star})$
\end{algorithmic}
\end{algorithm}

\subsection{Training and Its Limitations}
Training uses $168$ hand-authored synthetic utterances (the module's own
docstring labels them ``not logged real user queries''), a seeded $80/20$
split ($134$ train, $34$ validation), Adam ($\text{lr}=0.01$, weight decay
$10^{-3}$), cross-entropy loss, $150$ epochs, and best-validation-accuracy
early stopping. The saved checkpoint reports \textbf{training accuracy
$1.00$ and validation accuracy $0.529$} --- a textbook overfitting
signature on a tiny synthetic set~\cite{goodfellow2016deep}. An independent
re-evaluation reproduced exactly $0.529$ and produced the per-class
breakdown in Section~\ref{sec:results}: four of the twelve classes have
$F_1 = 0$. \textbf{We therefore make no generalization claim for this
classifier.} Its role is strictly a fallback signal behind the
deterministic rules, and even in that role it is gated at $P \ge 0.6$.
The $12$ fine-grained classes (e.g.\ separating \code{ROOM\_LOCATION} from
\code{DEPARTMENT\_LOCATION}) collapse onto the $5$ agents; the finer labels
exist so the classifier reports a useful intent even though several map to
the same agent.

\subsection{Design Rationale}
The deterministic-first design is justified by four constraints from
Section~\ref{sec:problem}: (i)~zero routing latency matters when the
generation step already costs seconds; (ii)~the deterministic rules are
zero-false-negative for the phrasings they cover and were not tuned to pass
any particular test; (iii)~every routing decision must be explainable from
logs; and (iv)~no labeled routing corpus exists to train a primary
classifier on. The LSTM adds coverage for paraphrases outside the keyword
lists (``What can I study here?'' contains no course keyword) without
making routing confidence-dependent for the common case.
